\chapter{Sistema adaptativo y concept drift}
\label{cap:adaptacion}

El benchmark anterior sirve como línea base no adaptativa. Sobre él se implementó una evaluación por
bloques temporales: cada bloque futuro se predice con las transacciones disponibles hasta ese momento,
usando una ventana expansiva o una ventana deslizante de tamaño fijo. El protocolo permite observar la
degradación del modelo, el efecto de olvidar historia y el costo de las decisiones.

\section{Features causales}

Se conserva la idea del \emph{magic UID} del capítulo anterior, pero sus estadísticas se calculan de
forma causal. Para una transacción en el instante $t$, la frecuencia y los promedios del UID usan solo
observaciones anteriores a $t$. En el bloque de evaluación se emplean mapas construidos con el histórico
previo; ninguna fila del futuro actualiza sus propias features antes de ser predicha. Esta corrección
elimina la limitación transductiva identificada en
\cref{sec:limites}.

También se derivan el día relativo, la hora, el día de la semana, el mes calendario y versiones
normalizadas de las columnas $D$. El identificador se elimina antes del entrenamiento: solo funciona
como llave de agregación.

\section{Ventanas y detección de drift}

El experimento utiliza 60 días iniciales y bloques futuros de 30 días. Se comparan cuatro
estrategias: histórico completo (ventana expansiva) y ventanas deslizantes de 30, 60 y 90 días.
Para cada bloque se calcula el Population Stability Index (PSI) sobre monto, variables $D$ y
conteos $C$. Se marca drift distributivo cuando el PSI máximo supera 0,20; el rendimiento también
se monitorea mediante AUC, PR-AUC, F1, recall, balanced accuracy y costo de decisión.

\begin{figure}[ht]
  \centering
  \includegraphics[width=0.98\textwidth]{\figdir fig_arquitectura_adaptativa.pdf}
  \caption{Arquitectura propuesta: el monitoreo temporal alimenta la selección de ventana y el
           reentrenamiento del modelo.}
  \label{fig:arquitectura-adaptativa}
\end{figure}

\section{Resultados de la evaluación adaptativa}

La corrida completa usa las 590\,540 transacciones de entrenamiento, XGBoost con 180 árboles,
umbral de decisión 0,20 y costos $1$ para un falso positivo y $10$ para un falso negativo. Se
reservan los bloques de los días 60--90, 90--120, 120--150 y 150--180 para evaluación futura.

\begin{figure}[ht]
  \centering
  \includegraphics[width=0.99\textwidth]{\figdir fig_adaptacion_metricas.pdf}
  \caption{Evolución de AUC, F1 y PSI por bloque. El descenso del segundo bloque muestra que el
           rendimiento puede degradarse aunque el drift distributivo no supere el umbral alto de
           PSI; por ello se monitorean también métricas de desempeño.}
  \label{fig:adaptacion-metricas}
\end{figure}

La ventana expansiva obtuvo un AUC promedio de $0,9163$ y un F1 promedio de $0,5575$. Las ventanas
de 60 y 90 días alcanzaron AUC promedio de $0,9139$ y $0,9156$, mientras que la ventana de 30 días
obtuvo $0,9039$. En esta corrida, conservar más historia funcionó mejor: la ventana corta perdió
señal y no compensó el drift moderado observado. El AUC expansivo bajó de $0,9175$ en el primer bloque
a $0,9071$ en el segundo, y el costo de decisión subió de 18\,059 a 21\,009 unidades. Por eso el
sistema monitorea el desempeño aunque el PSI no active una alarma fuerte.

\section{El costo de la causalidad, medido en el holdout}
\label{sec:costo-causalidad}

\begin{sloppypar}
La brecha transductiva señalada en \cref{sec:limites} quedó cuantificada con el script
\texttt{src/eval\_05\_integrated.py}, que evalúa el sistema integrado del notebook 05 en dos variantes
y registra cada una en su CSV. La variante \emph{protocolo} replica el notebook y añade métricas de
decisión: reproduce sus AUC ($0{,}9340/0{,}9456$ en holdout y $0{,}9347/0{,}9509$ en OOF, sin y con
magic UID) y mide recall, precisión y F1 en los umbrales 0,50 y 0,19, además del umbral
$\tau^{*}=0{,}176$ que maximiza F1 en el OOF ($0{,}6631$, in-sample del umbral).
\end{sloppypar}

La variante \emph{causal} separa tres fases: entrenamiento en $[0\%,60\%)$, calibración del umbral en
$[60\%,75\%)$ y evaluación en el futuro $[75\%,100\%]$. Categorías, frecuencias y agregaciones del UID
se calculan solo con transacciones anteriores a cada corte. El AUC del holdout baja de $0{,}9456$ a
$0{,}9175$ ($-0{,}028$): la distancia entre el límite superior transductivo y lo que un despliegue
estricto obtendría.

\begin{table}[ht]
  \centering
  \small
  \caption{Sistema integrado 05 + magic UID, holdout causal 75/25: métricas de decisión por umbral
           (\texttt{results/threshold\_metrics\_05\_causal.csv}). El umbral $\tau^{*}$ se calibró por
           F1 en el bloque $[60\%,75\%)$ sin usar información futura.}
  \label{tab:causal-umbral}
  \begin{tabular}{lrrr}
    \toprule
    Umbral & Recall & Precisión & F1\\
    \midrule
    \rowcolor{mist}
    0,50 & 0,3218 & 0,9308 & 0,4782\\
    0,19 & 0,4200 & 0,8903 & 0,5707\\
    \rowcolor{mist}
    $\tau^{*}=0{,}108$ & \textbf{0,4602} & 0,8651 & \textbf{0,6008}\\
    \bottomrule
  \end{tabular}
\end{table}

El \cref{tab:causal-umbral} muestra el compromiso operativo. Con umbral fijo 0,19, la variante causal
recupera menos fraude que la transductiva (recall $0{,}4200$ frente a $0{,}5586$ en el holdout
comparable); la calibración recupera parte de esa pérdida sin mirar el futuro y sostiene F1 $0{,}60$
en el bloque posterior. La conclusión operativa es doble: el AUC transductivo es un límite superior, y
el umbral debe recalibrarse en cada bloque con datos pasados.

\section{Decisión y adaptación operativa}

El score se transforma en una acción con tres estados: aprobar, enviar a revisión o bloquear. El
umbral puede modificarse según la utilidad del negocio. En producción, el sistema conservaría el
modelo vigente mientras monitorea AUC y PSI. Ante una caída de rendimiento, un PSI alto o el cierre de
un nuevo bloque temporal, recalcularía las features causales y reentrenaría con la ventana seleccionada.
Todas las predicciones, decisiones y métricas deben quedar registradas para permitir auditoría humana
y rollback del modelo.
